\documentclass[10pt,a4paper]{amsart}
\usepackage[margin=19mm]{geometry}
\usepackage{amsmath,amssymb,mathtools}
\usepackage[scr=rsfs,cal=euler]{mathalfa}
\usepackage{xfrac}
\usepackage[unicode,hidelinks,bookmarks=false]{hyperref}
\newcommand{\ie}{i.e.\ }
\newcommand{\cf}{cf.\ }
\newcommand{\resp}{respectively\ }
\newcommand{\Subsection}{\subsection}
\providecommand{\nobreakdash}{\nobreak\hbox{-}}
\setlength{\emergencystretch}{2em}
\multlinegap0pt
\newtheorem{theorem}{Theorem}
\newtheorem{corollary}{Corollary}
\title{What induces plane structures in complete graph drawings?}
\author{Alexandra Weinberger, Ji Zeng}
\date{Source: arXiv:2603.05208v1 (CC BY 4.0)}
\begin{document}
\maketitle

\noindent\emph{Source and licence.} What induces plane structures in complete graph drawings?. Version arXiv:2603.05208v1 (CC BY 4.0), distributed by arXiv under the Creative Commons Attribution 4.0 International licence, http://creativecommons.org/licenses/by/4.0/, as recorded in the arXiv licence metadata for this exact version and shown on the abstract page https://arxiv.org/abs/2603.05208v1. Full licence notice: \texttt{licenses/arxiv-2603.05208-CC-BY-4.0.txt}.\par\smallskip

\noindent\textit{Editorial scope.} Counted unit: the article's corollary cor:main0 (source0.tex lines 81--86). The article's complete original proof of it is delivered with it (source0.tex lines 248--290, one \texttt{proof} environment of 43 lines, which the article places away from the statement). No mathematics is written, repaired or re-derived: the statement and the proof are the article's own lines, in the article's own notation, and no retained line is substituted or re-flowed.  Retained uncounted as the source-local context the proof needs: lines 50--52.  One declared adaptation: exactly 4 ancillary strings inside the retained spans are omitted (image-inclusion commands and, where the source repeats a label, the repeated label definitions) and no other line differs from the source; the figure environments, with their placement, captions and labels, are kept, and the package records the omitted strings in fidelity receipts bound to the affected bodies.  No retained line contains a citation, so the wrapper carries no bibliography and no bibliography entry.  The retained lines contain the article's own diagram code, which the wrapper typesets with the standard tikz packages; no image file is delivered and the package declares no asset dependency.  Editorial formatting only, in addition: an AMS \texttt{amsart} preamble that restates the article's own declarations and macros used by the retained lines, namely the environments \texttt{theorem}, \texttt{corollary} on separate counters, together with the standard packages those lines need.  The retained environments print in this document as Theorem 1, Corollary 1 in order of appearance, whereas the article prints the same items with the numbering of its own full document.  The complete native source of the article as distributed in the arXiv e-print for this exact version is bundled unchanged as \texttt{sources/195785/source0.tex}.
\begin{theorem}\label{thm:main1}
    For any positive integer $m$, there exists $n$ such that every complete graph drawing on $n$ vertices that is either adjacent-simple or separate-simple must contain $m$ pairwise disjoint edges. On the other hand, there exist complete graph drawings on arbitrarily many vertices in which any two adjacent edges cross exactly once, and any two non-adjacent edges cross at least once and at most twice.
\end{theorem}

\begin{corollary}\label{cor:main0}
    For any positive integer $m$, there exists $n$ such that the following is true. Let $p_1,p_2,\dots,p_n$ be distinct points in the plane, $c_{ij}: [0,1] \to \mathbb{R}^2$ be piecewise smooth function with $c_{ij}(0)=p_i$ and $c_{ij}(1)=p_j$ for all $1\leq i<j \leq n$, and suppose the family $\{c_{ij}\}$ is stroke-like. If either of the following conditions holds: \begin{enumerate}
        \item[(A)] $\forall$ $i<j$ and $k<\ell$ with $|\{i,j,k,\ell\}| = 3$, we have $|c_{ij}^{-1}(\text{Im}(c_{k\ell}))| = |c_{k\ell}^{-1}(\text{Im}(c_{ij}))| = 1$;
        \item[(S)] $\forall$ $i<j$ and $k<\ell$ with $|\{i,j,k,\ell\}| = 4$, we have $|c_{ij}^{-1}(\text{Im}(c_{k\ell}))| = |c_{k\ell}^{-1}(\text{Im}(c_{ij}))| \leq 1$, and there is no touching point between $c_{ij}$ and $c_{k\ell}$;
    \end{enumerate} then there exist $m$ functions in $\{c_{ij}:~1\leq i<j \leq n\}$ whose images are pairwise disjoint.
\end{corollary}

\begin{proof}[Proof of Corollary~\ref{cor:main0}]
By abuse of notation, we shall call each $p_i$ a vertex and each $c_{ij}$ (or its image) an edge. We shall perform local modifications to the given configuration of vertices and edges, and obtain a drawing with mild assumptions that is either adjacent-simple or separate-simple. We make sure not to introduce any new pair of disjoint edges, so that Theorem~\ref{thm:main1} can be applied to conclude the proof.

Condition (A) means any two adjacent edges have exactly one intersection counted via pre-images. In particular, an edge $ab$ cannot pass through another vertex $c$ (otherwise $ac$ and $ab$ have two intersections), and an edge $ab$ cannot pass through $a$ or $b$ in its interior (otherwise another edge incident to $a$ or $b$ intersects $ab$ twice counted via pre-images). We now eliminate degeneracies without introducing disjointness or violating adjacent-simplicity.

\begin{figure}[ht]
    \centering
    
    \caption{We do not perturb a segment (such as $\alpha_2$) sandwiched by others.}
    \label{fig:perturb_touch}
\end{figure}

\begin{figure}[ht]
    \centering
    
    \caption{A strand is locally twisted to avoid introducing disjointness.}
    \label{fig:unpack}
\end{figure}

\begin{figure}[ht]
    \centering
    
    \caption{A standard local redraw to resolve self-crossings.}
    \label{fig:selfcross}
\end{figure}

\begin{figure}[ht]
    \centering
    
    \caption{Any two adjacent edges intersect before and after the perturbation.}
    \label{fig:perturb_block}
\end{figure}

First, if there is a touching point $p$ (whether a self-touching or between distinct edges), we locally perturb the drawing near $p$ to decrease the total number of touchings (counted via pre-images). A segment $\alpha_2$ is said to be \textit{sandwiched} if there are segments $\alpha_1$ and $\alpha_3$ that touch it in $p$ from two sides. To avoid introducing disjointness, we always perturb a segment through $p$ that is not sandwiched there, see Figure~\ref{fig:perturb_touch}. The existence of such a segment is justified at the end of the proof. Since touchings only happen between non-adjacent edges under Condition (A), the new crossings introduced in this step will not violate adjacent-simplicity.

Next, we remove coincidences between edges along segments. For any maximal segment $\sigma$ along which edges coincide, its multiplicity is defined as the total number of pre-image sub-intervals mapping onto $\sigma$ among all the functions $c_{ij}$. We unpack $\sigma$ into that many duplicated parallel strands lying very close to each other. Since coincidences along segments do not happen between adjacent edges under condition (A), we may reconnect the resulting parallel strands arbitrarily (in a local manner) to form edges. Moreover, we can locally twist some strands to avoid introducing disjointness, see Figure~\ref{fig:unpack}.

Then, we resolve crossing points that are self-intersections by a standard redraw technique, see Figure~\ref{fig:selfcross}. Although our mild assumption does not forbid multiple segments crossing at the same point, we can assume such a scenario does not happen after slight perturbations, which makes it more intuitive to resolve self-crossings. At the end, we obtain an adjacent-simple complete graph drawing satisfying the mild assumptions as wanted.

Condition (S) means any two non-adjacent edges have at most one intersection counted via pre-images, and that no such intersection is a touching point. Under this condition, an edge $ab$ cannot pass through two other vertices $c$ and $d$, otherwise $ab$ and $cd$ have two intersections. In particular, an edge $ab$ cannot pass through two other vertices $c$ and $d$, since then $ab$ and $cd$ would meet at least twice (at $c$ and $d$). Therefore, for $n$ sufficiently large we may greedily select as many vertices as we wish so that, among the selected vertices, no induced edge contains other selected vertices in its interior. We restrict our considerations to these selected vertices. Moreover, we can use local perturbations to forbid any edge $ab$ passing through its endpoint $a$ in its interior, see Figure~\ref{fig:perturb_block}. Now we can remove touchings, unpack coincident segments, and resolve self-crossings as above. Since Condition (S) forbids non-adjacent edges to touch, the removal of touchings will not make two non-adjacent edges cross more than once. At the end, we obtain a separate-simple complete graph drawing satisfying the mild assumptions as wanted.

Finally, we verify the existence of non-sandwiched segments. Let $p$ be a touching point on a segment $\alpha$, and consider a sufficiently small closed disk $D$ centered at $p$ whose boundary circle meets the drawing transversely and contains no other touching or crossing points. The two endpoints of $\alpha\cap \partial D$ split $\partial D$ into two arcs; we fix one and call it the left arc of $\alpha$. For any segment $\beta$ that touches $\alpha$ at $p$, we say that $\beta$ touches $\alpha$ from the left if both endpoints of $\beta\cap \partial D$ lie on the left arc of $\alpha$. If no segment touches $\alpha$ from the left, then $\alpha$ is not sandwiched. Otherwise, we consider all segments touching $\alpha$ from the left. Such a segment $\beta$ determines its own left arc on $\partial D$ (namely, the arc of $\partial D$ cut out by $\beta\cap \partial D$ that lies entirely within the left arc of $\alpha$). We define a partial order by declaring $\gamma \succ \beta$ if the left arc of $\gamma$ is contained in the left arc of $\beta$. A maximal element of this partially ordered set cannot be sandwiched, hence such a non-sandwiched segment always exists.
\end{proof}

\end{document}
